% !TeX program = xelatex
% UTF-8; compile twice with XeLaTeX. No external images or data files needed.
\documentclass[a4paper,landscape,10pt]{article}
\usepackage[margin=14mm,headheight=14pt,headsep=5mm,footskip=7mm]{geometry}
\usepackage{fontspec}
\setmainfont{DejaVu Sans}
\newfontfamily\ipafont{DejaVu Sans}
\usepackage[table]{xcolor}
\usepackage{array,longtable,booktabs,fancyhdr}
\definecolor{ink}{HTML}{18384D}
\definecolor{band}{HTML}{E7EFF4}
\pagestyle{fancy}\fancyhf{}
\fancyhead[L]{\small\color{ink}Parlez Vous QSO?! · Română / Romanian}
\fancyhead[R]{\small Englische Phrase · Traducere română · IPA}
\fancyfoot[L]{\scriptsize Vorlage: Romanian v 1.1 · Credits: Maria Szarvaczki}
\fancyfoot[R]{\small\thepage}
\setlength{\parindent}{0pt}
\setlength{\tabcolsep}{3mm}
\renewcommand{\arraystretch}{1.25}
\newcommand{\ipa}[1]{{\ipafont[#1]}}
\begin{document}
{\LARGE\bfseries\color{ink}Parlez Vous QSO?!}\hfill{\large Română / Romanian}

\vspace{2mm}
{\small QSO-Sprachführer mit rumänischer IPA-Lautschrift. Die Lautschrift bezieht sich auf die rumänische Spalte.}

\vspace{2mm}
{\footnotesize \textbf{Leseschlüssel:} \ipa{ˈ} Hauptbetonung, \ipa{ˌ} Nebenbetonung;
\ipa{ə} = ă, \ipa{ɨ} = â/î, \ipa{ʃ} = ș, \ipa{ts} = ț, \ipa{tʃ} = ce/ci, \ipa{dʒ} = ge/gi;
\ipa{ʲ} bezeichnet Palatalisierung, \ipa{e̯a, o̯a} Diphthonge. Grundlage: rumänische Standardsprache.
Schrägstriche trennen Alternativen; \ipa{…} steht für einen Platzhalter.}

\vspace{1mm}
\fontsize{9}{11.8}\selectfont
\begin{longtable}{@{}>{\raggedright\arraybackslash}p{13mm}>{\raggedright\arraybackslash}p{73mm}>{\raggedright\arraybackslash}p{75mm}>{\raggedright\arraybackslash}p{90mm}@{}}
\toprule
\textbf{Ref.} & \textbf{English phrase} & \textbf{Română} & \textbf{Lautschrift (IPA)}\\
\midrule
\endfirsthead
\toprule
\textbf{Ref.} & \textbf{English phrase} & \textbf{Română} & \textbf{Lautschrift (IPA)}\\
\midrule
\endhead
\midrule
\multicolumn{4}{r}{\scriptsize Fortsetzung auf der nächsten Seite}\\
\endfoot
\bottomrule
\endlastfoot
\rowcolor{band} 1.0 & \textbf{Opening Greetings} & \multicolumn{2}{l}{\textbf{Salutări}} \\*
1.1 & I am very pleased to meet you. & Mă bucur să vă întâlnesc.\newline Mă bucur să te întâlnesc. & \ipa{mə ˈbukur sə və ɨntɨlˈnesk / mə ˈbukur sə te ɨntɨlˈnesk} \\
1.2 & Thank you for your call & Mulțumesc pentru apel. & \ipa{multsuˈmesk ˈpentru aˈpel} \\
1.3 & Good morning & Bună dimineața. & \ipa{ˈbunə dimiˈne̯atsa} \\
1.31 & Good day & Bună ziua. & \ipa{ˈbunə ˈziwa} \\
1.32 & Good afternoon & Bună ziua. & \ipa{ˈbunə ˈziwa} \\
1.33 & Good evening & Bună seara. & \ipa{ˈbunə ˈse̯ara} \\
1.4 & I know only a little Romanian but I enjoy QSOs in Romanian. & Eu vorbesc doar puțin română,\newline dar îmi plac QSO-urile în română. & \ipa{jew vorˈbesk do̯ar puˈtsin roˈmɨnə dar ɨmʲ plak kju es o.urile ɨn roˈmɨnə} \\
\rowcolor{band} 2.0 & \textbf{Operator} & \multicolumn{2}{l}{\textbf{Operator}} \\*
2.1 & My name is Peter. & Mă numesc Peter. & \ipa{mə nuˈmesk ˈpiːtə} \\
2.2 & I have been licensed since 1960. & Am din 1960 licența pentru radioamatorism. & \ipa{am din o ˈmi.e ˈno.wə ˈsute ʃajˈzetʃʲ liˈtʃentsa ˈpentru radio.amatoˈrism} \\
2.3 & I have been licensed for 43 years. & Sunt de 43 ani radioamator. & \ipa{sunt de ˈpatruzetʃʲ ʃi trej anʲ radio.amaˈtor} \\
2.4 & I live in Sittingbourne. & Locuiesc în Sittingbourne. & \ipa{lokuˈjesk ɨn ˈsɪtɪŋbɔːn} \\
2.5 & Sittingbourne is a village. & Sittingbourne este un sat. & \ipa{ˈsɪtɪŋbɔːn ˈjeste un sat} \\
2.51 & Sittingbourne is a town. & Sittingbourne este o localitate. & \ipa{ˈsɪtɪŋbɔːn ˈjeste o lokaliˈtate} \\
2.52 & Sittingbourne is a city. & Sittingbourne este un oraș. & \ipa{ˈsɪtɪŋbɔːn ˈjeste un oˈraʃ} \\
2.6 & Sittingbourne is between London and Dover. & Sittingbourne este situat între Londra și Dover. & \ipa{ˈsɪtɪŋbɔːn ˈjeste situˈat ˈɨntre ˈlondra ʃi ˈdəʊvə} \\
2.7 & I live in the South of England. & Eu trăiesc în Anglia de Sud. & \ipa{jew trəˈjesk ɨn ˈaŋɡli.a de sud} \\
2.71 & I live in the North of England. & Eu trăiesc în Anglia de Nord. & \ipa{jew trəˈjesk ɨn ˈaŋɡli.a de nord} \\
2.72 & I live in the East of England. & Eu trăiesc în Anglia de Est. & \ipa{jew trəˈjesk ɨn ˈaŋɡli.a de est} \\
2.73 & I live in West of England. & Eu trăiesc în Anglia de Vest. & \ipa{jew trəˈjesk ɨn ˈaŋɡli.a de vest} \\
2.74 & I live in South East England. & Eu trăiesc în sud-estul Angliei. & \ipa{jew trəˈjesk ɨn ˌsudˈestul ˈaŋɡli.ej} \\
2.75 & I live in South West England. & Eu trăiesc în sud-vestul Angliei. & \ipa{jew trəˈjesk ɨn ˌsudˈvestul ˈaŋɡli.ej} \\
2.76 & I live in North East England. & Eu trăiesc în nord-estul Angliei. & \ipa{jew trəˈjesk ɨn ˌnordˈestul ˈaŋɡli.ej} \\
2.77 & I live in North West England. & Eu trăiesc în nord-vestul Angliei. & \ipa{jew trəˈjesk ɨn ˌnordˈvestul ˈaŋɡli.ej} \\
2.78 & I live in central England. & Eu trăiesc în Anglia centrală. & \ipa{jew trəˈjesk ɨn ˈaŋɡli.a tʃenˈtralə} \\
2.8 & I live in the country. & Eu locuiesc la țară. & \ipa{jew lokuˈjesk la ˈtsarə} \\
2.81 & I live by the sea. & Eu locuiesc în apropierea litoralului. & \ipa{jew lokuˈjesk ɨn apropiˈere̯a litoˈraluluj} \\
2.82 & I live near mountains. & Eu locuiesc în apropierea munților. & \ipa{jew lokuˈjesk ɨn apropiˈere̯a ˈmuntsilor} \\
2.9 & Sittingbourne is 50Km from London. & Sittingbourne este la 50 km de Londra. & \ipa{ˈsɪtɪŋbɔːn ˈjeste la tʃinˈzetʃʲ de kiloˈmetri de ˈlondra} \\
2.11 & I am retired. & Eu sunt pensionar. & \ipa{jew sunt pensi.oˈnar} \\
2.12 & My home call sign is M0ECS. & Indicativul meu de bază este M0ECS. & \ipa{indikaˈtivul mew de ˈbazə ˈjeste em ˈzero e tʃe es} \\
\rowcolor{band} 3.0 & \textbf{Station} & \multicolumn{2}{l}{\textbf{Stația}} \\*
3.1 & My radio is an Icom 706. & Aparatul meu de radioamatorism este Icom 706. & \ipa{apaˈratul mew de radio.amatoˈrism ˈjeste ˈaɪkɒm ˈʃapte ˈsute ˈʃase} \\
3.2 & My station is all homebrew (homemade). & Stația mea este completă și de mine construită. & \ipa{ˈstatsi.a me̯a ˈjeste komˈpletə ʃi de ˈmine konstruˈitə} \\
3.3 & I am running 5 Watts. & Eu transmit cu 5 wați. & \ipa{jew transˈmit ku tʃintʃʲ vatsʲ} \\
3.4 & My antenna is a long wire. & Antena mea este o sârmă lungă. & \ipa{anˈtena me̯a ˈjeste o ˈsɨrmə ˈluŋɡə} \\
3.41 & My antenna is a dipole. & Antena mea este un dipol. & \ipa{anˈtena me̯a ˈjeste un diˈpol} \\
3.42 & My antenna is a G5RV. & Antena mea este un G5RV. & \ipa{anˈtena me̯a ˈjeste un ɡe tʃintʃʲ er ve} \\
3.43 & My antenna is a doublet. & Antena mea este un Doublet. & \ipa{anˈtena me̯a ˈjeste un ˈdʌblɪt} \\
3.44 & My antenna is a delta loop. & Antena mea este un Delta-Loop. & \ipa{anˈtena me̯a ˈjeste un ˈdelta luːp} \\
3.45 & My antenna is a vertical. & Antena mea este verticală. & \ipa{anˈtena me̯a ˈjeste vertiˈkalə} \\
3.46 & My antenna is a log periodic. & Antena mea este un Log-Periodic. & \ipa{anˈtena me̯a ˈjeste un ˌlɒɡ pɪəriˈɒdɪk} \\
3.47 & My antenna is a Yagi beam. & Antena mea este un Yagi. & \ipa{anˈtena me̯a ˈjeste un ˈjaɡi} \\
3.48 & My antenna is a 5 element beam. & Antena mea este din 5 elemente. & \ipa{anˈtena me̯a ˈjeste din tʃintʃʲ eleˈmente} \\
\rowcolor{band} 4.0 & \textbf{Conditions} & \multicolumn{2}{l}{\textbf{Legăturile}} \\*
4.1 & Thank you for my 599 report. & Mulțumesc pentru raportul meu bun. & \ipa{multsuˈmesk ˈpentru raˈportul mew bun} \\
4.2 & The band conditions are good. & Legăturile sunt bune. & \ipa{leɡəˈturile sunt ˈbune} \\
4.3 & The band conditions are bad. & Legăturile sunt proaste. & \ipa{leɡəˈturile sunt ˈpro̯aste} \\
4.4 & There is strong interference here. & Am bruiaje puternice. & \ipa{am bruˈjaʒe puˈternitʃe} \\
4.5 & Can you change to the 80m band? & Puteți schimba pe banda de 80 metri? & \ipa{puˈtetsʲ skimˈba pe ˈbanda de ˈoptzetʃʲ ˈmetri} \\
4.6 & Please change frequency to 3.558 MHz. & Vă rog schimbați pe 3,558 MHz. & \ipa{və roɡ skimˈbatsʲ pe trej ˈvirɡulə tʃintʃʲ tʃintʃʲ opt ˈmeɡaˌherts} \\
\rowcolor{band} 5.0 & \textbf{General Queries and Comments} & \multicolumn{2}{l}{\textbf{Întrebări generale și comentarii}} \\*
5.1 & What does <\#\#\#\#\#> mean please? & Ce înseamnă, vă rog, <\#\#\#\#\#>? & \ipa{tʃe ɨnˈse̯amnə və roɡ …} \\
5.2 & Please wait a moment, I need to check my Romanian dictionary. & Așteptați un moment, mă uit în dicționarul român. & \ipa{aʃtepˈtatsʲ un moˈment mə ujt ɨn diktsi.oˈnarul roˈmɨn} \\
5.3 & How do you say, <\#\#\#\#\#> in Romanian? & Cum se spune <\#\#\#\#\#> în română? & \ipa{kum se ˈspune … ɨn roˈmɨnə} \\
5.5 & My favourite modes of operation are CW / RTTY / PSK / MT63 / Olivia / Hellschreiber / Digital Modes / and SSB & Forma mea de comunicare cea mai plăcută este: telegrafia / RTTY / PSK / MT63 / Olivia / Hell / forme digitale / și fonie. & \ipa{ˈforma me̯a de komuniˈkare tʃe̯a maj pləˈkutə ˈjeste teleɡraˈfi.a / er te te i ɡrek / pe es ka / em te ʃajˈzetʃʲ ʃi trej / oˈlivi.a / hɛl / ˈforme didʒiˈtale / ʃi foˈni.e} \\
5.6 & My best DX is Australia. & Legătura mea cea mai îndepărtată a fost în Australia. & \ipa{leɡəˈtura me̯a tʃe̯a maj ɨndepərˈtatə a fost ɨn a.usˈtrali.a} \\
5.7 & I enjoy working contests. & Particip cu plăcere la concursuri. & \ipa{partiˈtʃip ku pləˈtʃere la konˈkursurʲ} \\
5.8 & I like working QRP. & Lucrez cu drag cu parametrii mici / cu QRP. & \ipa{luˈkrez ku draɡ ku paraˈmetri mitʃʲ / ku kju er pe} \\
\rowcolor{band} 6.0 & \textbf{The Weather} & \multicolumn{2}{l}{\textbf{Vremea}} \\*
6.1 & It is sunny. & Este însorit. & \ipa{ˈjeste ɨnsoˈrit} \\
6.2 & It is windy. & Este vântos. & \ipa{ˈjeste vɨnˈtos} \\
6.3 & It is raining. & Plouă. & \ipa{ˈplo.wə} \\
6.4 & It is snowing. & Ninge. & \ipa{ˈnindʒe} \\
6.5 & It is very hot. & Este foarte cald. & \ipa{ˈjeste ˈfo̯arte kald} \\
6.6 & It is very cold. & Este foarte frig. & \ipa{ˈjeste ˈfo̯arte friɡ} \\
6.7 & The temperature is 8C. & Temperatura este de opt grade. & \ipa{temperaˈtura ˈjeste de opt ˈɡrade} \\
6.8 & It is overcast. & Este înnorat. & \ipa{ˈjeste ɨnoˈrat} \\
6.9 & It is foggy. & Este ceață. & \ipa{ˈjeste ˈtʃe̯atsə} \\
6.11 & It has been raining all day. & A plouat toată ziua. & \ipa{a ploˈwat ˈto̯atə ˈziwa} \\
6.12 & There is 20cm of snow here. & Sunt 20 cm zăpadă. & \ipa{sunt ˈdo.wəzetʃʲ de tʃentiˈmetri zəˈpadə} \\
\rowcolor{band} 8.0 & \textbf{QSL} & \multicolumn{2}{l}{\textbf{QSL / Cartela QSL}} \\*
8.1 & My QSL card is sure via the bureau & Cartela mea QSL vine prin birou. & \ipa{karˈtela me̯a kju es el ˈvine prin biˈrow} \\
8.2 & QSL via eQSL & Eu transmit un eQSL. & \ipa{jew transˈmit un i kju es el} \\
8.3 & No eQSL please & Nu doresc eQSL. & \ipa{nu doˈresk i kju es el} \\
8.4 & I will send you my card direct. & Eu transmit QSL-ul meu direct. & \ipa{jew transˈmit kju es elul mew diˈrekt} \\
8.5 & What is your address? & Care este adresa dumneavoastră? & \ipa{ˈkare ˈjeste aˈdresa dumne̯aˈvo̯astrə} \\
8.6 & Please send your QSL card direct. & Vă rog transmiteți QSL-ul direct. & \ipa{və roɡ transˈmitetsʲ kju es elul diˈrekt} \\
8.7 & My address is... & Adresa mea este… & \ipa{aˈdresa me̯a ˈjeste …} \\
8.8 & My address is shown on QSL.com / EHam.net & Adresa mea este trecută în QSL.com / eHam.net. & \ipa{aˈdresa me̯a ˈjeste treˈkutə ɨn kju es el punkt kom / i hæm punkt net} \\
8.9 & Sorry I do not send QSL cards. & Scuzați, nu transmit cartela QSL. & \ipa{skuˈzatsʲ nu transˈmit karˈtela kju es el} \\
8.11 & Please send QSL card via G1JQH. & Transmiteți, vă rog, QSL-ul prin G1JQH. & \ipa{transˈmitetsʲ və roɡ kju es elul prin ɡe unu ʒe kju haʃ} \\
8.12 & My WACRAL / GQRP / DXCC / FISTS number is 1234. & Numărul meu WACRAL / GQRP / DXCC / FISTS este 1234. & \ipa{ˈnumərul mew ˈwakral / ɡe kju er pe / de iks tʃe tʃe / fɪsts ˈjeste o ˈmi.e ˈdo.wə ˈsute ˈtrejzetʃʲ ʃi ˈpatru} \\
\rowcolor{band} 9.0 & \textbf{Final Greetings} & \multicolumn{2}{l}{\textbf{Salutări de final / Rămas-bun}} \\*
9.1 & Thank you. & Mulțumesc. & \ipa{multsuˈmesk} \\
9.2 & Thank you for a very enjoyable contact. & Mulțumesc de legătura frumoasă. & \ipa{multsuˈmesk de leɡəˈtura fruˈmo̯asə} \\
9.3 & I enjoyed our chat. & M-a bucurat a putea vorbi cu dv. & \ipa{ma bukuˈrat a puˈte̯a vorˈbi ku dumne̯aˈvo̯astrə} \\
9.4 & I enjoyed our QSO in Romanian. & M-am bucurat de QSO-ul în limba română. & \ipa{mam bukuˈrat de kju es oul ɨn ˈlimba roˈmɨnə} \\
9.5 & It has been great fun. & M-am bucurat.\newline Mi-a făcut plăcere. & \ipa{mam bukuˈrat / mja fəˈkut pləˈtʃere} \\
9.6 & If you would like a regular contact with me, you can contact me via email. & Dacă vreți să mă întâlniți regulat, trimiteți-mi un e-mail. & \ipa{ˈdakə vretsʲ sə mə ɨntɨlˈnitsʲ reɡuˈlat triˈmitetsʲmi un ˈiːmeɪl} \\
9.7 & My email address is... & Adresa mea de e-mail este… & \ipa{aˈdresa me̯a de ˈiːmeɪl ˈjeste …} \\
9.8 & My email address is shown on QSL.com / EHam.net & Adresa mea de e-mail este trecută în QSL.com / EHam.net. & \ipa{aˈdresa me̯a de ˈiːmeɪl ˈjeste treˈkutə ɨn kju es el punkt kom / i hæm punkt net} \\
9.9 & I look forward to meeting you again. & M-aș bucura să vă reîntâlnesc. & \ipa{maʃ bukuˈra sə və reɨntɨlˈnesk} \\
9.11 & My best wishes to you and your family. & Cele mai frumoase urări dv. și familiei. & \ipa{ˈtʃele maj fruˈmo̯ase uˈrərʲ dumne̯aˈvo̯astrə ʃi faˈmili.ej} \\
9.12 & Best wishes to your wife. & Cele mai frumoase urări soției. & \ipa{ˈtʃele maj fruˈmo̯ase uˈrərʲ soˈtsi.ej} \\
9.121 & Best wishes to your husband. & Cele mai frumoase urări soțului. & \ipa{ˈtʃele maj fruˈmo̯ase uˈrərʲ ˈsotsului} \\
9.13 & Good luck. & Noroc mult. & \ipa{noˈrok mult} \\
9.14 & Good health. & Sănătate bună. & \ipa{sənəˈtate ˈbunə} \\
9.15 & Goodbye & La reauzire. & \ipa{la re.a.uˈzire} \\
9.16 & I wish you an enjoyable afternoon. & Vă doresc o zi frumoasă. & \ipa{və doˈresk o zi fruˈmo̯asə} \\
9.161 & I wish you an enjoyable evening. & Vă doresc o seară plăcută. & \ipa{və doˈresk o ˈse̯arə pləˈkutə} \\
9.162 & I wish you an enjoyable weekend. & Vă doresc un sfârșit de săptămână frumos. & \ipa{və doˈresk un sfɨrˈʃit de səptəˈmɨnə fruˈmos} \\
9.17 & Enjoy your holiday. & Un concediu frumos. & \ipa{un konˈtʃedju fruˈmos} \\
9.18 & Happy New Year. & Un an bun. & \ipa{un an bun} \\
\rowcolor{band} 10.0 & \textbf{Country} & \multicolumn{2}{l}{\textbf{Țara}} \\*
10.1 & Albania & Albania & \ipa{alˈbani.a} \\
10.2 & Austria & Austria & \ipa{ˈa.ustri.a} \\
10.3 & Belgium & Belgia & \ipa{ˈbeldʒi.a} \\
10.4 & Belorussia & Belarus & \ipa{belaˈrus} \\
10.5 & Bosnia \& Herzegovina & Bosnia-Herțegovina & \ipa{ˈbosni.a hertseɡoˈvina} \\
10.6 & Bulgaria & Bulgaria & \ipa{bulˈɡari.a} \\
10.7 & Croatia & Croația & \ipa{kroˈatsi.a} \\
10.8 & Czech Republic & Cehia & \ipa{ˈtʃehi.a} \\
10.9 & Denmark & Danemarca & \ipa{daneˈmarka} \\
10.11 & Eire & Irlanda & \ipa{irˈlanda} \\
10.12 & England & Anglia & \ipa{ˈaŋɡli.a} \\
10.13 & Estonia & Estonia & \ipa{esˈtoni.a} \\
10.14 & Finland & Finlanda & \ipa{finˈlanda} \\
10.15 & France & Franța & \ipa{ˈfrantsa} \\
10.16 & Germany & Germania & \ipa{dʒerˈmani.a} \\
10.17 & Gibraltar & Gibraltar & \ipa{dʒibralˈtar} \\
10.18 & Greece & Grecia & \ipa{ˈɡretʃi.a} \\
10.19 & Holland & Olanda & \ipa{oˈlanda} \\
10.21 & Hungary & Ungaria & \ipa{uŋˈɡari.a} \\
10.22 & Iceland & Islanda & \ipa{isˈlanda} \\
10.23 & Northern Ireland & Irlanda de Nord & \ipa{irˈlanda de nord} \\
10.24 & Isle of Man & Insula Man & \ipa{ˈinsula man} \\
10.25 & Italy & Italia & \ipa{iˈtali.a} \\
10.26 & Latvia & Letonia & \ipa{leˈtoni.a} \\
10.27 & Liechtenstein & Liechtenstein & \ipa{ˈlɪçtn̩ʃtaɪn} \\
10.28 & Lithuania & Lituania & \ipa{lituˈani.a} \\
10.29 & Luxembourg & Luxemburg & \ipa{ˈluksemburɡ} \\
10.31 & Macedonia & Macedonia & \ipa{matʃeˈdoni.a} \\
10.32 & Moldavia & Moldova & \ipa{molˈdova} \\
10.33 & Netherlands & Olanda & \ipa{oˈlanda} \\
10.34 & Norway & Norvegia & \ipa{norˈvedʒi.a} \\
10.35 & Poland & Polonia & \ipa{poˈloni.a} \\
10.36 & Portugal & Portugalia & \ipa{portuˈɡali.a} \\
10.37 & Romania & România & \ipa{romɨˈni.a} \\
10.38 & Russian Federation & Rusia & \ipa{ˈrusi.a} \\
10.39 & Scotland & Scoția & \ipa{ˈskotsi.a} \\
10.41 & Slovakia & Slovacia & \ipa{sloˈvatʃi.a} \\
10.42 & Slovenia & Slovenia & \ipa{sloˈveni.a} \\
10.43 & Spain & Spania & \ipa{ˈspani.a} \\
10.44 & Svalbard & Svalbard & \ipa{ˈsvɑːlbɑr} \\
10.45 & Sweden & Suedia & \ipa{suˈedi.a} \\
10.46 & Switzerland & Elveția & \ipa{elˈvetsi.a} \\
10.47 & Turkey & Turcia & \ipa{ˈturtʃi.a} \\
10.48 & Ukraine & Ucraina & \ipa{uˈkrajna} \\
10.49 & Wales & Țara Galilor & \ipa{ˈtsara ˈɡalilor} \\
\end{longtable}
\normalsize
\section*{Hinweise zur Übertragung und Aussprache}
\fontsize{8.5}{10}\selectfont\raggedright
\textbf{Vorlage und Umfang.} Grundlage ist „Romanian v 1.1.doc“, Credit: Maria Szarvaczki.
Der innere Titel nennt abweichend v 1.0; hier gilt die Versionsangabe des Dateinamens.
Alle 142 Phrasen und Ländernamen sind enthalten. Die Gliederung und Nummernsprünge bleiben erhalten,
auch der fehlende Abschnitt 7. Referenz „132“ ist zu „1.32“ berichtigt.

\smallskip
\textbf{Schreibweise.} Rumänische Diakritika wurden wiederhergestellt: ă, â, î, ș und ț,
mit Komma unter ș/ț. Ältere î-Schreibungen im Wortinneren wurden auf die heutige â-Schreibung umgestellt,
wo erforderlich; reîntâlnesc behält î an der Wortgrenze. Offensichtliche Tippfehler wurden korrigiert,
etwa dimimeata, und oras, apropiereea, nea, inourat, Saluari und reantalnesc.
Abstände, Satzzeichen, Bindestriche und Zeilenumbrüche wurden vereinheitlicht.

\smallskip
\textbf{Inhaltliche Korrekturen.} „telefon“ wurde im Funkkontext zu „apel“ (1.2).
Die auf Englisch bezogenen Passagen 1.4, 5.2 und 9.4 wurden in beiden Sprachspalten auf Rumänisch umgestellt.
1.4 ist entsprechend dem englischen Sinn neu formuliert.
„Codul localității mele este MOECS“ wurde zu „Indicativul meu de bază este M0ECS“ (2.12);
„G5GV“ zu „G5RV“ (3.42), jeweils gemäß der englischen Spalte.
„Vați“ steht als „wați“ (3.3), „pe român“ als „în română“ (5.3),
„contestații“ als „concursuri“ (5.7, Wettbewerbe) und „La reauzite“ als „La reauzire“ (9.15).

\smallskip
\textbf{Ländernamen.} Die vier fehlenden Übersetzungen wurden ergänzt:
Belarus (10.4), Insula Man (10.24), Svalbard (10.44), Țara Galilor (10.49).
Bosnia-Herțegovina, Liechtenstein und Lituania ersetzen die fehlerhaften Schreibweisen der Vorlage.
Macedonia, Rusia und die doppelte Angabe Olanda für Holland/Netherlands bleiben erhalten.

\smallskip
\textbf{Besonderheiten der Vorlage.} 4.1 enthält im Rumänischen keinen 599-Rapport;
4.2/4.3 sprechen von Verbindungen statt Bandbedingungen. 5.5 verwendet telegrafia, Hell und fonie
statt wortwörtlich CW, Hellschreiber und SSB. 6.7 nennt acht Grad ohne Celsius;
6.11 steht in der Vergangenheit. 9.16 wünscht einen schönen Tag statt ausdrücklich einen Nachmittag.
Diese Unterschiede bleiben erhalten, ebenso einige ältere oder knappe Formulierungen.
Die IPA-Spalte folgt der rumänischen Spalte; „dv.“ wird als „dumneavoastră“ gelesen.

\smallskip
\textbf{Aussprachebasis.} Die manuelle, breite IPA-Umschrift orientiert sich an der rumänischen
Standardsprache. \ipa{ə} steht für ă, \ipa{ɨ} für â/î; \ipa{j, w} sind Gleitlaute.
\ipa{ʲ} bezeichnet ein kurzes palatalisiertes Ende, etwa bei ani und vielen Formen auf -ți;
bei metri und centimetri bleibt ein voller i-Vokal. Eu und este werden mit anlautendem \ipa{j} gelesen.
R ist vereinfacht als \ipa{r} dargestellt. Satzmelodie, alle Assimilationen und regionale Varianten
sind nicht vollständig kodiert. Cincizeci steht mit der gebräuchlichen Variante „cinzeci“.

\smallskip
\textbf{Namen und Kürzel.} Sittingbourne, Dover, Peter, Icom, Doublet, Loop und Log-Periodic
verwenden englische Aussprache; Hell und Liechtenstein deutsche, Svalbard norwegische.
Andere Namen sind überwiegend rumänisch gelesen. Rufzeichen und Kürzel verwenden hier Q „chiu“,
S „es“, L „el“, M „em“, R „er“, G „ghe“, C „ce“, J „je“, H „haș“, Y „i grec“.
Andere Buchstabennamen sind ebenfalls gebräuchlich. WACRAL und FISTS werden als Wörter gelesen.
Das e in eQSL/eHam ist \ipa{i}; Ham und e-mail stehen englisch, Webpunkte als „punct“, com/net als Wörter.

\smallskip
\textbf{Zahlen und Platzhalter.} 1960 = „o mie nouă sute șaizeci“, 43 = „patruzeci și trei“,
706 = „șapte sute șase“, MT63 mit „șaizeci și trei“, 1234 = „o mie două sute treizeci și patru“.
3,558 MHz wird als „trei virgulă cinci cinci opt megahertz“ gelesen; Maßeinheiten werden ausgeschrieben.
Zeitklammern bei Tagesgrüßen werden nicht mitgesprochen. Offene Angaben und \texttt{<\#\#\#\#\#>}
stehen als Auslassung.

\smallskip
\textbf{Fachhinweise.} Ioana Chitoran, Ioana Vasilescu, Bianca Vieru und Lori Lamel:
„Connected speech in Romanian: Exploring sound change through an ASR system“,
öffentliches Manuskript auf vocapia.com, besonders Tabellen 1/2 zum Lautsystem.
DEX/DOOM-Einträge auf dexonline.ro wurden für Schreibungen, Betonungen und Buchstabennamen herangezogen.
\end{document}
